\documentclass[11pt]{article}
\usepackage[a4paper,margin=28mm]{geometry}
\usepackage[T1]{fontenc}
\usepackage{lmodern}
\usepackage{microtype}
\usepackage{xcolor}
\usepackage{enumitem}
\usepackage[hidelinks]{hyperref}

\definecolor{revisionblue}{RGB}{0,0,255}
\setlength{\parindent}{0pt}
\setlength{\parskip}{0.7em}
\setlist[itemize]{leftmargin=1.6em,itemsep=0.35em,topsep=0.3em}

\begin{document}

\begin{center}
{\Large Revision of \emph{Construction of Kochen--Specker Sets from Mutually Unbiased Bases}}\\[0.5em]
{\small 4 August 2026}
\end{center}

Dear Mirko,

I have now worked carefully through both reports and carried out a detailed revision of the manuscript.  The first referee recommends publication without revision.  I found the second referee's requests reasonable: the paper benefited from a shorter introduction, earlier definitions, a clearer statement of the principal results, a guide to the larger construction, and a sharper separation between the higher-dimensional forcing arguments and the coordinate inventories.

I have therefore prepared a revised manuscript, with all substantive additions and changed passages marked in blue, together with separate draft responses to the referee and the editor.  I also checked the ray, context, and reconstruction inventories independently.  That audit revealed several points that go beyond the referee's presentational requests and that, in my view, had to be corrected before resubmission.

\textbf{Presentation and organization.}
\begin{itemize}
\item The Introduction is substantially shorter and now defines rays, contexts, intertwining rays, two-valued states, Kochen--Specker configurations, separating families, and unitality when they first become relevant.
\item The main claims are isolated as three propositions: the exact qutrit provenance inventory and the forcing constraints in dimensions four and five.
\item The qutrit construction begins with a four-stage guide, making clear the roles of the seeds, the 25 algebraic maps, projective identification, and exhaustive orthogonality enumeration.
\item In the higher-dimensional part, the forcing equations and their proofs now precede, and are separated from, the completion rays and computational inventories.
\item Repetition, speculative wording, and overlong figure captions have been reduced.  The distinction between a selected context hypergraph and the complete orthogonality hypergraph of the same ray set is now explicit throughout.
\end{itemize}

\textbf{Corrections resulting from the exact audit.}
\begin{itemize}
\item \emph{Qutrit provenance.}  The previously reported $40$--$4$--$4$ ``asymmetry'' was a labeling artifact.  Eighteen rays common to all three generated lineages had been assigned to the first lineage simply because its labels were encountered first.  Projectively invariant accounting gives 21 rays common to all three lineages and 48 rays exclusive to each lineage.  At the context level there are 10 common contexts and 40 lineage-specific contexts for each lineage, of which four use only exclusive rays.  Thus each lineage has 69 rays and 50 contexts, and the union has 165 rays and 130 contexts.  The revised tables and discussion now state this symmetric decomposition.

\item \emph{The 20-ray four-dimensional inventory.}  Exhaustive orthogonality enumeration gives 17 contexts, not 13: the 10 imposed contexts, three automatic disjoint-support contexts, and four equal-modulus contexts.  The selected 10- and 13-context hypergraphs each have 36 separating two-valued states, whereas the complete 17-context hypergraph has 16.  The text now also observes the necessary monotonicity: adding contexts can preserve or reduce the set of two-valued states, but cannot create new ones.

\item \emph{Comparison with the Cabello set.}  The gadget and Cabello configurations share 15 rays projectively, not 16.  In particular, $v_{1423}=(1,-1,-1,1)$ is not collinear with the Cabello ray $v_{47}=(1,1,-1,1)$.  I replaced that identification with two separate rank-three orthogonal-complement reconstructions and corrected three further reconstruction triples.  The final table shows five gadget-only rays and three Cabello-only rays, with mutual reconstruction from the other set.

\item \emph{The former four-dimensional MUB appendix.}  Each displayed collection was an orthonormal basis after normalization, but the five bases were not pairwise mutually unbiased.  For example, the computational vector $(1,0,0,0)$ and the displayed Bell vector $(1,0,0,1)/\sqrt{2}$ have squared overlap $1/2$, rather than $1/4$.  More decisively, the displayed $\mathcal B_1$ and $\mathcal B_4$ shared the normalized rays $(1,1,1,1)/2$ and $(1,-1,1,-1)/2$, giving squared overlap 1.  The displayed $\mathcal B_2$ was also essentially a $Y\otimes X$ basis rather than the claimed $Y\otimes Y$ basis.  Since the dimension-four proof needs only the equal-modulus condition $|x_1|^2=\cdots=|x_4|^2=1/4$, not a catalogue of five MUBs, I removed the incorrect catalogue and stated the required criterion directly.  The verified qutrit MUB appendix has been retained.
\end{itemize}

The revised source compiles cleanly, with resolved citations and cross-references, and I have checked the resulting ten-page PDF page by page.  The enclosed response drafts describe both the referee-requested changes and these author-initiated corrections transparently.

Could you please look especially closely at the four audit corrections above and let me know whether you agree with their formulation before we submit the revision?  Of course, please also edit the response letters as you see fit.

Best wishes,

Karl

\end{document}
